\documentclass[12pt]{article}
\usepackage[utf8]{inputenc}
\usepackage{geometry}
\usepackage{listings}
\usepackage{xcolor}
\usepackage{hyperref}

\geometry{margin=1in}

\definecolor{codebg}{RGB}{240,240,240}
\definecolor{keyword}{RGB}{0,0,180}
\definecolor{comment}{RGB}{0,128,0}

\lstdefinestyle{ruststyle}{
    backgroundcolor=\color{codebg},
    language=Rust,
    basicstyle=\ttfamily\small,
    keywordstyle=\color{keyword}\bfseries,
    commentstyle=\color{comment}\itshape,
    breaklines=true,
    frame=single,
    showstringspaces=false,
    tabsize=2
}

\title{Membrane Computation and Max-SMT in Rust}
\author{}
\date{\today}

\begin{document}

\maketitle

\tableofcontents
\newpage

\section{Membrane Computation Overview}

Membrane computation is a parallel computational framework where:

\begin{itemize}
    \item \textbf{Membranes} represent candidate solutions.
    \item \textbf{Rules} transform membrane states (variables) iteratively.
    \item \textbf{Environment} defines constraints.
    \item \textbf{Evaluation} determines which membranes survive.
\end{itemize}

\subsection{Binary SMT Approach}
Membranes are either \textbf{alive} if constraints are satisfied, or \textbf{dead} otherwise.

\subsection{Max-SMT / Optimization Approach}
Each membrane is assigned a \textbf{score} based on how many constraints it satisfies.
Soft constraints contribute partially, while hard constraints can still eliminate membranes.

\[
\text{score}(M) = \sum_{i=1}^{n} w_i \cdot \text{sat}(c_i, M)
\]

where \(w_i\) is the weight of constraint \(c_i\) and \(\text{sat}(c_i, M)\) is the degree of satisfaction (0..1).

\section{Rust Implementation}

\subsection{Membrane and Constraint Structures}

\begin{lstlisting}[style=ruststyle]
use serde::Serialize;
use std::collections::HashMap;

#[derive(Serialize)]
struct Membrane {
    id: usize,
    variables: HashMap<String, f64>,
    score: f64,
}

#[derive(Clone)]
struct Constraint {
    name: String,
    weight: f64,
    limit: f64,
    vars: Vec<String>,
}

impl Constraint {
    fn satisfaction(&self, vars: &HashMap<String, f64>) -> f64 {
        let sum: f64 = self.vars.iter().map(|v| *vars.get(v).unwrap_or(&0.0)).sum();
        if sum <= self.limit {
            1.0
        } else {
            self.limit / sum
        }
    }
}
\end{lstlisting}

\subsection{Evaluate Membrane Scores}

\begin{lstlisting}[style=ruststyle]
fn evaluate_membrane(m: &mut Membrane, constraints: &[Constraint]) {
    let mut total_score = 0.0;
    for c in constraints {
        total_score += c.weight * c.satisfaction(&m.variables);
    }
    m.score = total_score;
}
\end{lstlisting}

\subsection{Generate and Rank Membranes}

\begin{lstlisting}[style=ruststyle]
fn main() -> std::io::Result<()> {
    let constraints = vec![
        Constraint { name: "c1".to_string(), weight: 1.0, limit: 10.0, vars: vec!["x".to_string(), "y".to_string()] },
        Constraint { name: "c2".to_string(), weight: 2.0, limit: 5.0, vars: vec!["x".to_string()] },
    ];

    let mut membranes = vec![];
    for i in 0..5 {
        let mut vars = HashMap::new();
        vars.insert("x".to_string(), i as f64 + 1.0);
        vars.insert("y".to_string(), (5 - i) as f64);
        let mut m = Membrane { id: i, variables: vars, score: 0.0 };
        evaluate_membrane(&mut m, &constraints);
        membranes.push(m);
    }

    membranes.sort_by(|a, b| b.score.partial_cmp(&a.score).unwrap());

    let json = serde_json::to_string_pretty(&membranes)?;
    std::fs::write("membranes.json", json)?;
    println!("JSON saved to membranes.json");

    Ok(())
}
\end{lstlisting}

\subsection{Example D3.js JSON Output}

\begin{lstlisting}[style=ruststyle]
[
  {"id":0,"variables":{"x":1.0,"y":5.0},"score":3.9},
  {"id":1,"variables":{"x":2.0,"y":4.0},"score":4.2}
]
\end{lstlisting}

\section{Mutation Rules for Membranes}

\begin{lstlisting}[style=ruststyle]
use rand::Rng;

impl Membrane {
    fn mutate(&mut self, mutation_rate: f64, mutation_range: f64) {
        let mut rng = rand::thread_rng();
        for value in self.variables.values_mut() {
            if rng.gen::<f64>() < mutation_rate {
                let delta = rng.gen_range(-mutation_range..mutation_range);
                *value += delta;
            }
        }
    }
}
\end{lstlisting}

\subsection{Using Mutation in Evolution Loop}

\begin{lstlisting}[style=ruststyle]
for generation in 0..10 {
    for m in membranes.iter_mut() {
        evaluate_membrane(m, &constraints);
    }

    membranes.sort_by(|a, b| b.score.partial_cmp(&a.score).unwrap());

    for m in membranes.iter_mut().skip(1) {
        m.mutate(0.5, 1.0);
    }
}
\end{lstlisting}

\section{Visualization Notes}

\begin{itemize}
    \item Export membrane scores each generation to JSON.
    \item Color nodes by score (heatmap/gradient).
    \item Optional: animate over generations to see optimization progress.
\end{itemize}

\section{Summary}

\begin{itemize}
    \item Membrane computing provides a \textbf{parallel framework} for Max-SMT problems.
    \item Ranking membranes by score supports \textbf{soft constraints} and \textbf{optimization}.
    \item Mutation rules allow \textbf{exploration of solution space}.
    \item D3.js JSON output enables \textbf{interactive visualization}.
\end{itemize}

This framework can be extended to:
\begin{itemize}
    \item Non-linear SMT constraints
    \item Larger populations
    \item Advanced selection and mutation strategies
    \item Real-time interactive visualization
\end{itemize}

\end{document}
